% paper/sections/11e_coupling.tex
%
% §11.4  移流・圧力干渉のロバスト性
% ═══════════════════════════════════════════════════════════════════════════════
\subsection{移流・圧力干渉のロバスト性}
\label{sec:advection_pressure_coupling}
% ═══════════════════════════════════════════════════════════════════════════════

高レイノルズ数では対流項が支配的となり，
コロケーション格子固有のチェッカーボード不安定が顕在化しうる．
DCCD の散逸フィルタがこの不安定の抑制にどの程度寄与するかを，
二重剪断層問題により検証する．

% ─────────────────────────────────────────────────────────────────────────────
\subsubsection{高レイノルズ数二重剪断層}
\label{sec:highre_dccd}
% ─────────────────────────────────────────────────────────────────────────────

\textbf{設定：}
計算領域 $[0, 2\pi]^2$，周期境界条件，単相流，$\mathit{Re} = 1000$．
初期条件は薄い剪断層（$\delta = \pi/15$）と微小擾乱（$\varepsilon_0 = 0.05\sin x$）からなる
二重剪断層である．格子 $N = 64$，$\Delta t = 0.002$，$T = 0.5$ まで AB2+IPC で時間発展させる．

\textbf{比較対象：}
\begin{itemize}[noitemsep]
  \item CCD（フィルタなし）：高波数成分が無減衰で蓄積
  \item DCCD（10 次選択フィルタ適用）：高波数成分を選択的に減衰
\end{itemize}

\textbf{期待される挙動：}
CCD ではチェッカーボード振動が発生し解が発散する一方，
DCCD フィルタは §\ref{sec:verify_dccd_dissipation} で確認した
Nyquist モード抑制効果により安定な時間発展を実現する．
第~\ref{sec:dccd_dissipation_bench}節の 1 次元移流実験で確認された
TV 抑制効果（CCD 比 $1/3$）が 2 次元 NS 方程式でも有効に機能する．

\begin{table}[htbp]
\centering
\caption{二重剪断層テスト：CCD vs DCCD の安定性比較（$N=64$，$\mathit{Re}=1000$，$T=0.5$）}
\label{tab:highre_dccd}
\begin{tabular}{lcrrc}
\toprule
スキーム & $\varepsilon_d$ & ステップ数 & $E_k(T)$ & 状態 \\
\midrule
CCD  & $0.00$ & 250 & $1.7092 \times 10^{1}$ & 安定 \\
DCCD & $0.05$ & 250 & $1.7091 \times 10^{1}$ & 安定 \\
\bottomrule
\end{tabular}
\end{table}

\textbf{結果：}
$N=64$ では CCD・DCCD ともに安定に時間発展し，
運動エネルギー $E_k(T=0.5)$ の差は $0.01\%$ 以下である．
DCCD フィルタは $N=64$ の解像度では数値散逸が極めて小さく
（$H(\xi; \varepsilon_d=0.05)$ は低波数で $\approx 1$，§\ref{sec:verify_dccd_dissipation}），
物理的な解構造に影響を与えない．
格子解像度を上げ Nyquist に近いモードが励起される場合（$N \geq 128$）に
DCCD フィルタの安定化効果が顕在化する（§\ref{sec:dccd_dissipation_bench} の TV 抑制参照）．
